\section{Relation to Process Tensors, Causal Models, Contextuality, and Measurement}
\label{sec:relations}

This work does not propose new dynamics or a replacement formalism for quantum theory. It instead
provides a structural diagnostic: it identifies when many-to-one operational reductions are
incapable—by construction—of supporting counterfactual reasoning under a specified intervention
class, and it characterizes the minimal representational structure required to restore such
reasoning.

\subsection{Process tensors and quantum combs}

Multi-time process tensors and quantum combs provide operational descriptions of non-Markovian
processes by encoding multi-time correlations into algebraic objects. In the present classification,
these are refined $\Sigma_1$ descriptions: they extend beyond instantaneous states but remain
\emph{algebraically closed}, in the sense that prediction proceeds by contraction or composition
rules internal to a fixed tensorial object.

The structural result of Section~3 applies equally to such descriptions. If the operational cut
identifying a process object is degenerate with respect to correlations that can later be probed
by interventions, then static algebraic closure does not guarantee faithfulness. In such cases,
distinct admissible histories may correspond to the same process tensor while supporting different
counterfactual continuations.

\paragraph{Schematic failure mode.}
Consider two admissible histories $h_1,h_2\in\mathcal{H}$ that induce the same static process tensor
$\mathcal{T}$ over a fixed set of times, so that all observable statistics for insertions at those
times coincide. Suppose, however, that $h_1$ and $h_2$ differ in correlations with degrees of freedom
discarded by the reduction defining $\mathcal{T}$. If the intervention class $\mathcal{I}$ includes
a later operation (e.g., free evolution or a control pulse) whose effect depends on those discarded
correlations, then the same $\mathcal{T}$ supports distinct admissible future outcomes depending on
which history was realized.

In this situation, the process tensor correctly predicts outcomes for the interventions it was
constructed to represent, but it cannot support counterfactual prediction for the enlarged
intervention class. The failure is not dynamical but representational: the tensor is a many-to-one
image of $\mathcal{H}$ that discards continuation-relevant distinctions.

From the $\Sigma_2$ perspective, the appropriate completion is not a single enlarged tensor, but a
family of algebraic descriptions indexed by fiber elements or continuation-equivalence classes in
$\mathcal{F}(r)$. Equivalently, one may view this as a sheaf-like structure in which $\Sigma_1$
descriptions are valid only on patches of the fiber determined by the intervention context.

This perspective is compatible with existing results emphasizing that extension theorems for
stochastic or quantum processes can fail in the presence of active interventions. Here, that
failure is diagnosed as an instance of an unfaithful operational cut rather than as a limitation
of a particular tensor formalism.

\subsection{Quantum causal models and interventions}

Quantum causal models and process-matrix frameworks focus on admissibility constraints at the level
of operations and correlations, often seeking conditions under which observed statistics admit a
causal explanation. The present framework addresses a prior question: whether the reduced
descriptions assigned to nodes support intervention semantics at all.

If nodal reductions are $\Sigma_1$ and unfaithful with respect to the intervention class relevant to
a causal analysis, then no causal structure built atop those reduced objects can consistently
support counterfactual reasoning without incorporating fiber-level admissibility structure. In
this sense, $\Sigma_2$ supplies a diagnostic precondition for causal modeling in reduced operational
spaces, rather than a competing causal formalism.

\subsection{Contextuality as unfaithful reduction}

Preparation contextuality provides a canonical instance of an unfaithful cut: distinct preparation
procedures yield operationally equivalent reduced descriptions (the same $\rho$) yet differ in
their statistics under subsequent measurement interventions. In the present language, $\Pi$ is the
map from preparation histories to reduced states, and unfaithfulness is precisely the failure of
continuation invariance over the fiber $\Pi^{-1}(\rho)$.

From the $\Sigma_2$ perspective, the appropriate descriptive object for preparations is not $\rho$
alone, but the lifted structure $\Sigma_2(\rho)$, which retains the intervention-indexed
continuation geometry required to support counterfactual prediction. This reframes contextuality
as a geometric feature of the reduction rather than as a purely interpretive puzzle about what the
quantum state ``really is.''

\subsection{Measurement and event structure}

Measurement problems can likewise be viewed through the lens of unfaithful cuts. Operational
reductions that identify macroscopically distinct branch histories under a single reduced
description destroy the continuation distinctions relevant to counterfactual questions of the
form: ``what would occur if we measured observable $M$ rather than $N$?''

In this perspective, the issue is not primarily to find a modified reduced dynamics, but to
recognize that an unfaithful reduction cannot support event-level counterfactuals without
additional structure encoding branch- or history-level admissibility. $\Sigma_2$ provides a
minimal language for representing such structure without committing to a particular interpretive
ontology or modification of quantum dynamics.